\documentclass[11pt,a4paper]{article}
\usepackage[T1]{fontenc}
\usepackage[utf8]{inputenc}
\usepackage[brazil]{babel}
\usepackage[margin=2.2cm]{geometry}
\usepackage{booktabs}
\usepackage{enumitem}
\usepackage{fancyhdr}
\usepackage{hyperref}
\usepackage{setspace}

\hypersetup{colorlinks=true, linkcolor=black, urlcolor=blue}
\pagestyle{fancy}
\fancyhf{}
\fancyhead[L]{MC714 --- Trabalho 2}
\fancyhead[R]{Roteiro do vídeo (teleprompter)}
\fancyfoot[C]{\thepage}
\renewcommand{\headrulewidth}{0.4pt}
\onehalfspacing

\newcommand{\cue}[1]{\par\medskip\noindent\textit{[#1]}\par\medskip}

\begin{document}

\begin{center}
{\LARGE\bfseries Roteiro do vídeo --- Trabalho 2}\\[0.35em]
{\large MC714 --- Sistemas Distribuídos}\\[0.5em]
{\normalsize Rodrigo Camargo (238257) \quad André Loubet (253333) \quad Gabriel Eiji Tomonari (234774)}\\[0.3em]
{\small $\approx$ 5 minutos \quad / \quad $\approx$ 1:40 por pessoa}
\end{center}

\section*{Antes de gravar}
Cluster no ar, comandos copiados, fonte grande no teleprompter.

\section{0:00--1:40 --- Rodrigo}

\begin{quote}\large
Olá. Somos Rodrigo Camargo, André Loubet e Gabriel Eiji Tomonari. Este é o Trabalho 2 de MC714.

A gente implementou dois algoritmos principais. O primeiro é exclusão mútua, com Ricart-Agrawala. O segundo é um registro replicado ABD, com vários leitores e escritores.

Também juntamos os dois num contador. Para incrementar, o nó pega a seção crítica, lê, escreve valor mais um, e só então libera o lock.

Como bônus, tem um detector de falhas por heartbeat, reintegração depois de crash, e injeção de atraso e perda na rede.

Tudo isso roda em Go, por HTTP, cada nó num contêiner. RA precisa da resposta de todo mundo. ABD precisa só da maioria. Essa diferença aparece bem quando a gente mata um nó.
\end{quote}

\cue{docker compose ps}

\section{1:40--3:20 --- André}

\begin{quote}\large
Beleza. Cenário normal.

Três nós no ar. Vou pedir seções críticas concorrentes no Ricart-Agrawala. O observador só anota entrada e saída. Ele não dá permissão. A gente espera zero sobreposição.
\end{quote}

\cue{python3 experiments/ra\_smoke.py}

\begin{quote}\large
Sem falha, cada aquisição manda duas vezes N menos um mensagens.

Agora o ABD. Escrevo num nó e leio em outro. A escrita pergunta pra maioria, escolhe uma tag maior e espalha o valor. A leitura também pergunta pra maioria e faz write-back se precisar.
\end{quote}

\cue{write 42 no node1; read no node2}

\begin{quote}\large
No contador integrado, dez incrementos ao mesmo tempo deveriam virar de um a dez, com valor final dez. Nesse modo a escrita ABD direta fica bloqueada.
\end{quote}

\section{3:20--5:00 --- Gabriel}

\begin{quote}\large
Agora a falha.

Vou derrubar o node3. O detector de falhas deve marcar ele como suspeito. No RA, novas aquisições não concluem. No ABD, a maioria continua lendo e escrevendo.

O contador integrado depende do RA, então ele também para. A leitura ABD sozinha continua ok.
\end{quote}

\cue{kill node3; lab -action fd; tentar RA; read ABD}

\begin{quote}\large
Se a gente sobe o node3 de novo com allow-rejoin, ele sincroniza o registro pela maioria e volta a participar.

Na campanha grande, com três, cinco e sete nós, RA ficou em dez pedidos por segundo e ABD em vinte. Em partição, ABD seguiu na maioria; RA parou. Depois da reconexão, os dois voltaram.

Um cuidado: tudo isso foi numa máquina só. Contêiner não é máquina independente.

Código, testes, dados e relatório estão no repositório. Obrigado.
\end{quote}

\cue{tela final: repo + nomes}

\section*{Divisão}
\begin{center}
\begin{tabular}{@{}ll@{}}
\toprule
Rodrigo & intro e modelo \\
André & demo normal \\
Gabriel & falha, rejoin, fechamento \\
\bottomrule
\end{tabular}
\end{center}

\subsection*{Comandos}
\begin{verbatim}
python3 experiments/ra_smoke.py
./bin/lab -action write -url http://127.0.0.1:18081 -value 42 -wait-ms 3000
./bin/lab -action read  -url http://127.0.0.1:18082 -wait-ms 3000
docker compose kill node3
./bin/lab -action fd -url http://127.0.0.1:18081
./bin/lab -action ra -url http://127.0.0.1:18081 -hold-ms 50 -wait-ms 1000 -timeout 2s
./bin/lab -action read -url http://127.0.0.1:18081 -wait-ms 3000
\end{verbatim}

\end{document}
